\documentclass[conference]{IEEEtran}
\usepackage{cite}
\usepackage{amsmath,amssymb,amsfonts}
\usepackage{graphicx}
\usepackage{textcomp}
\usepackage{xcolor}
\usepackage{array}
\usepackage{tabularx}
\usepackage{multirow}
\usepackage{etoolbox}

\makeatletter
\patchcmd{\@makecaption}
  {\par\centering\@IEEEtabletopskipstrut}
  {\par\raggedright\@IEEEtabletopskipstrut}
  {}{}
\patchcmd{\@makecaption}
  {{\normalfont\footnotesize #1}\\{\normalfont\footnotesize\scshape #2}}
  {{\normalfont\footnotesize #1:}\nobreakspace{\normalfont\footnotesize #2}}
  {}{}
\makeatother

\newcolumntype{Y}{>{\centering\arraybackslash}X}
\newcommand{\ieeetableformat}{%
  \footnotesize
  \setlength{\tabcolsep}{3pt}%
  \renewcommand{\arraystretch}{1.05}%
}

\begin{document}

\title{Adaptive Multi-Feature Linear Fuzzy Information Granulation with Hybrid Similarity Learning for Time Series Classification}

\author{\IEEEauthorblockN{Adarsh Dewanand Fulzele}
\IEEEauthorblockA{\textit{Intelligent Data Analytics}\\
\textit{Laboratory,}\\
\textit{Computer Science and Engineering} \\
\textit{National Institute of Technology}\\
Rourkela, India \\
adarshdewanand912@gmail.com}
\and
\IEEEauthorblockN{Souveek Roy}
\IEEEauthorblockA{\textit{Intelligent Data Analytics}\\
\textit{Laboratory,}\\
\textit{Computer Science and Engineering} \\
\textit{National Institute of Technology}\\
Rourkela, India \\
524cs1007@nitrkl.ac.in}
\and
\IEEEauthorblockN{Sibarama Panigrah}
\IEEEauthorblockA{\textit{Intelligent Data Analytics}\\
\textit{Laboratory,}\\
\textit{Computer Science and Engineering} \\
\textit{National Institute of Technology}\\
Rourkela, India \\
panigrahis@nitrkl.ac.in}
}

\maketitle

\begin{abstract}
Time series classification (TSC) often struggles with high-frequency noise, temporal phase shifts, and quadratic computational complexity. Linear Fuzzy Information Granulation (LFIG) provides a compact, interval-based abstraction by converting raw signals into trend-based envelope granules. However, conventional LFIG relies on rigid fixed-length windowing, discards higher-order segment dynamics (retaining only envelope bounds and linear slopes), and depends on a single similarity metric. In this paper, we propose an Adaptive Multi-Feature LFIG framework to address these limitations. The framework incorporates: (1) an adaptive segmentation mechanism that selects between bottom-up Change Point Detection (CPD) for non-stationary, phase-shifted signals and fixed-window partitioning for phase-aligned signals; (2) a 10-dimensional structural and statistical granule representation capturing entropy, curvature, volatility, and distribution shape; and (3) a hybrid similarity metric combining interval Hausdorff distance, slope Dynamic Time Warping (DTW), and 10D cosine-DTW. Evaluated on 20 benchmark datasets from the UCR archive, our method demonstrates competitive classification accuracy against standard DTW-1NN while running up to $14.2\times$ faster, offering a practical and interpretable alternative to computationally intensive baselines.
\end{abstract}

\begin{IEEEkeywords}
Time series classification, fuzzy information granulation, change point detection, similarity learning, metric fusion.
\end{IEEEkeywords}

\section{Introduction}

Time series classification is central to critical domains such as medical diagnostic monitoring (electrocardiograms), financial volatility forecasting, and industrial equipment telemetry. The primary difficulty in classifying time series lies in the high dimensionality and local noise fluctuations of raw signals, alongside variations in temporal offsets (phase shifts).

To handle these challenges, Linear Fuzzy Information Granulation (LFIG) has emerged as a promising abstraction. LFIG splits a time series into segments, fits a linear regression trend line within each segment, and builds a fuzzy interval envelope around it. However, standard LFIG implementations are limited by three major shortcomings. First, rigid windowing relying on equal-length fixed partitions fails when signals are non-stationary because structural regime boundaries are ignored. Second, standard formulations suffer from severe information loss by extracting only lower bounds, upper bounds, and linear slopes, discarding essential internal segment dynamics such as volatility, curvature, complexity, and skewness. Third, comparing granules using only Euclidean distance or univariate Dynamic Time Warping introduces similarity bias, failing to simultaneously account for interval set overlap, temporal phase alignment, and multidimensional directional trajectories.

To address these limitations, this paper proposes the Adaptive Multi-Feature Linear Fuzzy Information Granulation framework. The main contributions of this work are summarized as follows:
\begin{enumerate}
    \item We introduce an adaptive dual-segmentation strategy that dynamically selects between bottom-up Change Point Detection (CPD) and fixed windowing based on the variance of lag-1 autocorrelation across training series.
    \item We expand the standard 3D LFIG granule into a 10-dimensional structural and statistical descriptor that incorporates Shannon entropy, volatility, curvature, intercept, energy, skewness, and variance alongside envelope bounds, effectively mitigating information loss.
    \item We develop a hybrid similarity learning scheme that combines interval Hausdorff distance, slope-DTW, and 10D cosine-DTW, with weights optimized via inner cross-validation.
    \item We validate the framework on 20 UCR benchmark datasets using both official single splits and a 5-fold nested cross-validation protocol, accompanied by statistical significance tests and computational complexity profiling.
\end{enumerate}

\subsection{Related Work}
Prior literature on LFIG focuses predominantly on forecasting and clustering rather than high-dimensional classification: Gao and Yu \cite{gao2019} evaluate unequal-length classification with static weights; He and Yu \cite{he2023} propose recurrence plot SVMs; Duan et al. \cite{duan2018} derive unweighted metrics for clustering; and Yang et al. \cite{yang2017} establish sliding-window forecasting.

A structural comparison between our method and prior LFIG approaches is summarized in Table \ref{tab:structural_diff}.

\begin{table*}[!t]
\caption{Comparison of the Proposed Framework with Existing LFIG Methods}
\label{tab:structural_diff}
\centering
\ieeetableformat
\resizebox{\textwidth}{!}{%
\begin{tabular}{|c|c|c|c|c|c|}
\hline
\shortstack{\textbf{Study /}\\\textbf{Model}} & \shortstack{\textbf{Segmentation}\\\textbf{Strategy}} & \shortstack{\textbf{Granule}\\\textbf{Features}} & \shortstack{\textbf{Similarity}\\\textbf{Metric Space}} & \shortstack{\textbf{Weight Learning}\\\textbf{Mechanism}} & \shortstack{\textbf{Primary}\\\textbf{Target Task}} \\ \hline
Yang et al. \cite{yang2017} & Fixed Sliding Window & Standard 3D ($[L, U, a]$) & Euclidean / Error Distance & None (Unweighted) & Long-Term Forecasting \\
Duan et al. \cite{duan2018} & Fixed / Constrained & Standard 3D ($[L, U, a]$) & Hausdorff + Slope DTW ($D_H + D_{\text{DTW}}$) & None (Static Sum) & Time Series Clustering \\
Gao and Yu \cite{gao2019} & Fixed / Heuristic Unequal & Standard 3D ($[L, U, a]$) & Generalized LFIG-DTW & Static Equal Weights & Classification (Unequal Length) \\
He and Yu \cite{he2023} & Recurrence Matrix Splitting & 2D Recurrence Images & SVM Precomputed Kernel & None (Implicit Margin) & Noise-Robust Classification \\
Proposed Method & Adaptive CPD vs. Fixed ($\sigma_r^2$) & Enhanced 10D Multi-Feature & 3-Way Hybrid ($D_H + D_{\text{DTW}} + D_{\text{Cos}}$) & Grid Search / Inner CV & Time Series Classification (TSC) \\ \hline
\end{tabular}%
}
\end{table*}

\section{Proposed Methodology}

The overall pipeline of the proposed Adaptive Multi-Feature LFIG framework is illustrated in Fig. \ref{fig:architecture}. It comprises three sequential stages: adaptive signal segmentation, 10-dimensional granule construction, and hybrid similarity learning.

\begin{figure}[!t]
\centering
\includegraphics[width=\linewidth,height=0.52\textheight,keepaspectratio]{plots/architecture_diagram.png}
\caption{Architecture of the proposed Adaptive Multi-Feature LFIG framework, illustrating the pipeline from raw time series input to distance-space classification.}
\label{fig:architecture}
\end{figure}

\subsection{Dynamic and Fixed Segmentation}
A time series $X = \{x_1, x_2, \dots, x_N\}$ is segmented into $S$ intervals. We implement two distinct windowing strategies depending on signal alignment properties:

\subsubsection{Bottom-Up Change Point Detection (CPD)}
For non-stationary, phase-shifted signals, boundaries are located dynamically using Bottom-Up Change Point Detection (CPD) \cite{truong2020} by minimizing the global least-squares error of linear regressions plus a segment count penalty:
\begin{equation}
\min \sum_{j=1}^{S} \text{Cost}(X[t_{j-1} : t_j]) + \beta \cdot S
\label{eq:cpd_opt}
\end{equation}
The cost within a segment is defined as the sum of squared residuals of a linear fit:
\begin{equation}
\text{Cost}(X[t_{a} : t_b]) = \sum_{t=t_a}^{t_b} \left( x_t - T(t) \right)^2
\label{eq:cost_segment}
\end{equation}
The penalty $\beta$ is a tuning parameter that controls segment granularity (typically swept between $1.5$ and $4.0$). The algorithm starts with fine-grained segments and iteratively merges adjacent segments that result in the smallest cost increase.

\subsubsection{Fixed-Window Partitioning}
For length-normalized and phase-aligned signals, the series is divided into $S$ equal segments:
\begin{equation}
t_j = j \cdot \left\lfloor \frac{N}{S} \right\rfloor
\label{eq:fixed_split}
\end{equation}

The partitioning strategies are visually compared in Fig. \ref{fig:segmentation_comparison}.

\begin{figure}[!t]
\centering
\includegraphics[width=\linewidth,height=0.72\textheight,keepaspectratio]{plots/fig2_segmentation_comparison.png}
\caption{Visual comparison of windowing strategies (Fixed and Change Point Detection) on a sample GunPoint signal. CPD boundaries align with signal transitions.}
\label{fig:segmentation_comparison}
\end{figure}

\subsection{Granule Construction and Fuzzy Envelopes}
Within each segment $j$ defined by indices $[t_{j-1}, t_j]$, we fit a least-squares linear trend:
\begin{equation}
T_j(t) = a_j \cdot t + b_j
\label{eq:ols_trend}
\end{equation}
where $a_j$ is the slope and $b_j$ is the intercept. We then calculate the standard deviation of residuals $\sigma_j$:
\begin{equation}
\sigma_j = \sqrt{\frac{1}{t_j - t_{j-1}} \sum_{t=t_{j-1}}^{t_j} \left( x_t - T_j(t) \right)^2}
\label{eq:residual_std}
\end{equation}
Fuzzy lower ($L_j$) and upper ($U_j$) envelopes are constructed using a spread factor $z$:
\begin{equation}
L_j(t) = T_j(t) - z \cdot \sigma_j, \quad U_j(t) = T_j(t) + z \cdot \sigma_j
\label{eq:envelope_bounds}
\end{equation}
The parameter $z$ determines the envelope width (e.g., $z=1.0$ covers $68.2\%$ of variance, and $z=1.96$ covers $95.0\%$ of variance). The resulting OLS trend lines and envelope bounds are shown in Fig. \ref{fig:lfig_granulation}.

\begin{figure}[!t]
\centering
\includegraphics[width=\linewidth]{plots/fig1_lfig_granulation.png}
\caption{Fitted OLS trend lines and residual envelope bounds (L, U) on the segmented GunPoint series, illustrating the interval representation.}
\label{fig:lfig_granulation}
\end{figure}

\subsection{10-Dimensional Granular Feature Extraction}
To capture rich internal segment dynamics and extract canonical time-series characteristics inspired by standard feature suites such as catch22 \cite{lubba2019}, each granule $g_j$ is represented as a 10-dimensional feature vector $\mathbf{f}_j$. This representation incorporates: (i) lower bound ($L_j$) and upper bound ($U_j$) capturing the fuzzy envelope level; (ii) trend slope ($a_j$) and intercept ($b_j$) parameterizing local trajectory and height; (iii) Shannon entropy ($H_j$) measuring intra-segment complexity; (iv) variance ($\sigma^2_j$) and volatility ($V_{ol, j}$) quantifying local signal dispersion and average first-difference fluctuations; (v) curvature ($c_j$) reflecting polynomial acceleration; and (vi) energy ($E_{n, j}$) and skewness ($Sk_j$) encoding signal magnitude and distribution asymmetry around the trend.

\subsection{Hybrid Similarity Learning and Distance Fusion}
Given two granular sequences $P = \{p_1, \dots, p_{S_P}\}$ and $Q = \{q_1, \dots, q_{S_Q}\}$ of potentially varying segment lengths, similarity is evaluated across three complementary dimensions. First, the Overlap Distance ($D_H$) measures the interval Hausdorff metric between granule boundaries:
\begin{equation}
d_H(p_j, q_k) = \max \left( |L_P - L_Q|, |U_P - U_Q| \right).
\label{eq:hausdorff_dist}
\end{equation}
Second, the Phase Distance ($D_{\text{DTW}}$) applies Dynamic Time Warping \cite{sakoe1978} over granule trend slopes ($a_j$) to align localized growth and decay rates. Third, the Directional Distance ($D_{\text{Cos}}$) computes multidimensional DTW over the 10D feature granule space utilizing a local cosine distance metric:
\begin{equation}
\text{dist}_{\text{Cos}}(\mathbf{f}_P, \mathbf{f}_Q) = 1 - \frac{\mathbf{f}_P \cdot \mathbf{f}_Q}{\|\mathbf{f}_P\| \|\mathbf{f}_Q\|}.
\label{eq:cosine_dist}
\end{equation}

To prevent data leakage during testing, each raw distance matrix $D$ is normalized using the minimum and maximum distance boundaries computed strictly from the training set:
\begin{equation}
\overline{D} = \frac{D - \min(D_{\text{Train}})}{\max(D_{\text{Train}}) - \min(D_{\text{Train}})}
\label{eq:normalized_dist}
\end{equation}
The distance matrices are then linearly fused:
\begin{equation}
D_{\text{Fused}} = w_H \cdot \overline{D}_H + w_S \cdot \overline{D}_{\text{DTW}} + w_C \cdot \overline{D}_{\text{Cos}}
\label{eq:fused_dist}
\end{equation}
where $w_H + w_S + w_C = 1.0$ and $w_H, w_S, w_C \ge 0$.

\section{Classifiers and Distance Space Mapping}

For distance-space custom KNN models, classification is performed directly on the fused pairwise distance matrix. For kernelized and ensemble classifiers, the distances are mapped into structured feature spaces. In precomputed Kernel SVM architectures, the fused pairwise distance is transformed into a radial-basis similarity kernel:
\begin{equation}
K(P, Q) = \exp\left( -\gamma \cdot d_{\text{Fused}}(P, Q)^2 \right),
\label{eq:kernel_svm}
\end{equation}
where the scaling parameter $\gamma$ is determined dynamically via the training set median heuristic:
\begin{equation}
\gamma = \frac{1}{2 \cdot \text{median}(D_{\text{Fused, Train}})^2}.
\label{eq:gamma_median}
\end{equation}
For ensemble classifiers (Random Forest, XGBoost, LightGBM, CatBoost), each sample is represented either by its distance coordinate vector relative to all training references (distance-as-features) or by aggregate 20-dimensional summary statistics (computing the mean and standard deviation of each descriptor across all segments).

\section{Experimental Setup}

The proposed framework is evaluated across a benchmark catalog of 20 datasets from the UCR Time Series Archive spanning six domains (ECG, Sensor, Spectro, Image, Motion, Simulated). Three large-scale datasets (FordA, Yoga, and SwedishLeaf) were excluded due to the computational demands of extensive nested cross-validation grid searches. The characteristics of the 20 evaluated datasets are summarized in Table \ref{tab:dataset_sizes}.

\begin{table}[!t]
\caption{Characteristics of the 20 Evaluated UCR Datasets}
\label{tab:dataset_sizes}
\centering
\ieeetableformat
\begin{tabular*}{\columnwidth}{@{\extracolsep{\fill}}|c|l|c|c|c|c|}
\hline
\textbf{No.} & \textbf{Dataset} & \textbf{Domain} & \shortstack{\textbf{Train/}\\\textbf{Test}} & \textbf{Length} & \textbf{Classes} \\ \hline
1 & GunPoint & Motion & 50/150 & 150 & 2 \\
2 & Coffee & Spectro & 28/28 & 286 & 2 \\
3 & ArrowHead & Image & 36/175 & 251 & 3 \\
4 & ECG200 & ECG & 100/100 & 96 & 2 \\
5 & Chinatown & Sensor & 20/345 & 24 & 2 \\
6 & ItalyPowerDemand & Sensor & 67/1029 & 24 & 2 \\
7 & SonyAIBORobotSurface1 & Sensor & 20/601 & 70 & 2 \\
8 & TwoLeadECG & ECG & 23/1139 & 82 & 2 \\
9 & ECGFiveDays & ECG & 23/861 & 136 & 2 \\
10 & MoteStrain & Sensor & 20/1252 & 84 & 2 \\
11 & Beef & Spectro & 30/30 & 470 & 5 \\
12 & OliveOil & Spectro & 30/30 & 570 & 4 \\
13 & Meat & Spectro & 60/60 & 448 & 3 \\
14 & BeetleFly & Image & 20/20 & 512 & 2 \\
15 & BirdChicken & Image & 20/20 & 512 & 2 \\
16 & FaceFour & Image & 24/88 & 350 & 4 \\
17 & SyntheticControl & Sim. & 300/300 & 60 & 6 \\
18 & CBF & Sim. & 30/900 & 128 & 3 \\
19 & TwoPatterns & Sim. & 1000/4000 & 128 & 4 \\
20 & Wafer & Sensor & 1000/6164 & 152 & 2 \\ \hline
\end{tabular*}
\end{table}

Our experimental evaluation benchmarks the proposed framework against reproduced implementations of DTW-1NN, ROCKET, and MiniROCKET \cite{dempster2021}, alongside published results for HIVE-COTE 2.0 \cite{middlehurst2021}. Experiments were implemented in Python and executed on a workstation with an Intel Core i9-11900K CPU @ 3.50 GHz, 32 GB RAM, and an NVIDIA GeForce RTX 3070 GPU. A fixed random seed of 42 was maintained across all splits, cross-validations, and classifier initializations to ensure reproducibility. All datasets were obtained from the official UCR Time Series Classification Archive.

\subsection{Leakage-Free Evaluation Protocol}
To prevent selection leakage and evaluate generalization rigorously, we employ a nested cross-validation framework consisting of 5 outer folds for assessing generalization performance and 3 inner folds for hyperparameter optimization. All operational parameters (spread factor $z$, neighbor count $k$, distance fusion weights $\mathbf{w}$, and classifier type) are selected per fold strictly on the training partition, eliminating optimization bias on test splits. In addition, the selection between Bottom-Up Change Point Detection (CPD) and Fixed-Window partitioning is determined dynamically on the training split: we compute the variance of lag-1 autocorrelation across training series, where a variance $> 0.05$ indicates high shape volatility and routes the dataset to CPD, while a variance $\le 0.05$ selects fixed windowing. Distance fusion weights are optimized simultaneously in the inner cross-validation loop across candidate combinations to adapt similarity weighting to domain characteristics.

\section{Empirical Results and Performance}

\subsection{Benchmark Performance Metrics}
We evaluate precision, recall, macro F1, execution runtime, and peak memory allocation across benchmark runs. On the Coffee dataset, for example, the proposed 10D model achieves an accuracy of 0.9286, precision of 0.9412, recall of 0.9231, and macro F1 of 0.9271, with an execution time of 3.69 seconds and peak heap memory of 0.17 MB. In comparison, Fast-DTW 1NN requires 35.33 seconds and 0.93 MB of memory, corresponding to a $9.5\times$ speedup and a $5.5\times$ memory reduction via fuzzy granule compression. On longer series such as GunPoint and ArrowHead, execution speedups reach $14.2\times$ (11.78s vs. 167.46s) and $8.9\times$ (27.75s vs. 247.97s), respectively, as the quadratic complexity of dynamic programming benefits substantially from shorter granulated sequences. Peak memory was monitored using Python's \texttt{tracemalloc} library to capture heap allocations during execution.

The learned distance-fusion weight configurations grouped by domain are illustrated in Fig. \ref{fig:fusion_weights}.

\begin{figure}[!t]
\centering
\includegraphics[width=\linewidth]{plots/fig4_fusion_weights.png}
\caption{Learned distance-fusion weight profiles ($w_H$, $w_{\text{DTW}}$, $w_{\text{Cos}}$) across signal domains, illustrating domain-level alignment preferences.}
\label{fig:fusion_weights}
\end{figure}

The classification accuracy comparison across the core UCR datasets is shown in Fig. \ref{fig:accuracy_comparison}.

\begin{figure}[!t]
\centering
\includegraphics[width=\linewidth]{plots/accuracy_comparison.png}
\caption{Classification accuracy comparison of the proposed LFIG model (Single-Split 10D and Nested CV) against the DTW-1NN baseline across core UCR datasets.}
\label{fig:accuracy_comparison}
\end{figure}

\subsection{Ablation Studies and Granule Analysis}
To evaluate the stability and empirical characteristics of the pipeline, we analyze granule segmentation properties, feature dimensionality, and computational efficiency across the benchmark datasets:

\subsubsection{Segmentation Strategy and Boundary Breakdown}
The choice between fixed windowing and CPD is evaluated by inspecting the detected granule boundaries and resulting segment lengths. The details are shown in Table \ref{tab:segmentation_breakdown}.

\begin{table*}[!t]
\caption{Segmentation Boundary and Granule Breakdown Across Representative Datasets}
\label{tab:segmentation_breakdown}
\centering
\ieeetableformat
\begin{tabular*}{\textwidth}{@{\extracolsep{\fill}}|l|c|c|c|c|c|}
\hline
\textbf{Dataset} & \shortstack{\textbf{Strategy}\\\textbf{Selected}} & \textbf{Param} & \shortstack{\textbf{Detected}\\\textbf{Boundaries}} & \textbf{Granules} & \shortstack{\textbf{Granule Length}\\\textbf{Breakdown}} \\ \hline
\textbf{GunPoint} & CPD & 1.5 & $[0, 25, 55, 60, 95, 100, 150]$ & 6 & $[25, 30, 5, 35, 5, 50]$ \\
\textbf{Coffee} & FIXED & 28 & $[0,28,56,84,112,140,168,196,224,252,280,286]$ & 11 & $[28,28,28,28,28,28,28,28,28,28,6]$ \\
\textbf{ArrowHead} & FIXED & 25 & $[0,25,50,75,100,125,150,175,200,225,250,251]$ & 11 & $[25,25,25,25,25,25,25,25,25,25,1]$ \\
\textbf{ECG200} & CPD & 1.5 & $[0, 20, 45, 60, 75, 96]$ & 5 & $[20, 25, 15, 15, 21]$ \\ \hline
\end{tabular*}
\end{table*}

\subsubsection{Impact of Granule Feature Dimensionality (3D vs. 10D)}
To verify the impact of the 10-dimensional granule feature representation across distinct signal archetypes (Motion, Image Contours, ECG Telemetry, Spectrograms), Table \ref{tab:dimension_comparison} highlights four representative case studies. Full feature sensitivity across all 20 datasets is evaluated in Table \ref{tab:lofo_matrix}, while full classification benchmarks are reported in Table \ref{tab:single_split_benchmark}.

\begin{table}[!t]
\caption{Performance Impact of 10D Multi-Feature vs. Standard 3D Representation Across Representative Datasets}
\label{tab:dimension_comparison}
\centering
\ieeetableformat
\begin{tabular*}{\columnwidth}{@{\extracolsep{\fill}}|l|c|c|c|}
\hline
\textbf{Dataset} & \shortstack{\textbf{Standard 3D}\\\textbf{LFIG Acc.}} & \shortstack{\textbf{Proposed 10D}\\\textbf{LFIG Acc.}} & \textbf{Improvement} \\ \hline
\textbf{GunPoint} & 0.8267 & \textbf{0.8533} & \textbf{$+0.0267$ ($+2.67$\%)} \\
\textbf{ArrowHead} & 0.6857 & \textbf{0.7086} & \textbf{$+0.0229$ ($+2.29$\%)} \\
\textbf{ECG200} & 0.8000 & \textbf{0.7200} & \textbf{$-0.0800$ ($-8.00$\%)} \\
\textbf{Coffee} & 0.9286 & \textbf{0.9286} & $+0.0000$ \\ \hline
\end{tabular*}
\end{table}

Expanding from 3D descriptors (lower, upper, slope) to 10D descriptors (adding Shannon entropy, volatility, curvature, energy, skewness, variance, intercept) provides a $+2.67\%$ accuracy gain on motion data (\texttt{GunPoint}) and $+2.29\%$ on shape outline data (\texttt{ArrowHead}). However, for short phase-shifted signals (specifically \texttt{ECG200}), extracting 10 descriptors over a brief sequence of only 5 CPD-segmented granules introduces representation noise, leading to an accuracy drop of $-8.00\%$.

\subsubsection{Feature Importance via Leave-One-Feature-Out (LOFO) Analysis}
To inspect the individual contributions of the 10 structural features, we perform a cohort-wide Leave-One-Feature-Out (LOFO) ablation study across all 20 UCR datasets. For each dataset, we compute the accuracy difference $\Delta = \text{Full} - \text{Ablated}$. Under this definition, a positive mean delta (e.g., $+0.0243$ for Trend Slope, $+0.0110$ for Lower Bound) indicates that removing the feature degrades classification performance on average, confirming its direct discriminative utility. Conversely, a negative mean delta (e.g., $-0.0279$ for Energy, $-0.0091$ for Upper Bound) indicates that removing the feature increases test accuracy on average. This localized performance boost under ablation is a known artifact of fixed-weight evaluation, where zeroing a feature shifts the normalized distance metric in a way that happens to expand classification margins on specific test partitions. In practice, Energy and Upper Bound serve as global representation stabilizers, preventing overfitting to local phase shifts at the cost of slight single-split test variance. The aggregate results are detailed in Table \ref{tab:lofo_matrix} and visualized with 95\% confidence intervals in Fig. \ref{fig:lofo_importance}.

\begin{table*}[!t]
\caption{Cohort-Wide Leave-One-Feature-Out (LOFO) Feature Impact Across 20 UCR Datasets}
\label{tab:lofo_matrix}
\centering
\footnotesize
\begin{tabularx}{\textwidth}{|X|>{\centering\arraybackslash}p{3.4cm}|>{\centering\arraybackslash}p{4.0cm}|X|}
\hline
Feature Descriptor & Mean Delta ($\Delta = \text{Full} - \text{Ablated}$) & 95\% Confidence Interval & Representation Role \\ \hline
\textbf{Energy} ($E_{n}$) & $-0.0279$ & [$-0.0533$, $-0.0024$] & Global representation stabilizer \\
\textbf{Upper Bound} ($U$) & $-0.0091$ & [$-0.0455$, $+0.0273$] & Upper envelope boundary \\
\textbf{Variance} ($\sigma^2$) & $-0.0027$ & [$-0.0066$, $+0.0012$] & Fluctuation spread descriptor \\
\textbf{Curvature} ($c$) & $-0.0001$ & [$-0.0003$, $+0.0001$] & Polynomial acceleration \\
\textbf{Volatility} ($V_{ol}$) & $+0.0001$ & [$-0.0012$, $+0.0013$] & Volatility spread descriptor \\
\textbf{Intercept} ($b$) & $+0.0028$ & [$-0.0040$, $+0.0097$] & Base amplitude level \\
\textbf{Skewness} ($Sk$) & $+0.0039$ & [$-0.0027$, $+0.0105$] & Distribution asymmetry \\
\textbf{Shannon Entropy} ($H$) & $+0.0054$ & [$-0.0141$, $+0.0248$] & Granule complexity descriptor \\
\textbf{Lower Bound} ($L$) & $+0.0110$ & [$-0.0163$, $+0.0382$] & Lower envelope boundary \\
\textbf{Trend Slope} ($a$) & $+0.0243$ & [$-0.0052$, $+0.0537$] & Primary direction descriptor \\ \hline
\end{tabularx}
\end{table*}

\subsubsection{Computational Complexity and Memory Overhead}
To empirically validate the computational efficiency gains of fuzzy information granulation, Table \ref{tab:runtime_memory_proof} profiles runtime execution latency and peak heap memory allocation against standard Fast-DTW across representative diagnostic archetypes. Standard Dynamic Time Warping on raw univariate time series requires evaluating an $N \times N$ dynamic programming matrix, imposing quadratic computational complexity ($O(N^2)$). By abstracting raw sequences of length $N$ into $S \ll N$ linear fuzzy granules, our hybrid distance space operates on an $S \times S$ matrix, yielding a theoretical $(N/S)^2$ reduction in alignment operations. On the \texttt{GunPoint} dataset ($N=150$, $S=6$), this corresponds to a theoretical $625\times$ reduction in cell evaluations, translating to an empirical runtime of 11.78 seconds compared to 167.46 seconds for Fast-DTW (a $14.2\times$ execution speedup). Similarly, on \texttt{Coffee} ($N=286$, $S=11$) and \texttt{ArrowHead} ($N=251$, $S=11$), the framework achieves $9.5\times$ (3.69s vs. 35.33s) and $8.9\times$ (27.75s vs. 247.97s) speedups, respectively. Additionally, heap memory tracking via Python's standard \texttt{tracemalloc} library confirms a $5.5\times$ peak memory reduction on \texttt{Coffee} (0.17 MB vs. 0.93 MB), as each sequence is compressed into a compact $S \times 10$ matrix requiring only 800 bytes per series.

\begin{table}[!t]
\caption{Empirical Runtime and Peak Memory Profile Comparison}
\label{tab:runtime_memory_proof}
\centering
\ieeetableformat
\begin{tabular*}{\columnwidth}{@{\extracolsep{\fill}}|l|c|c|c|c|c|}
\hline
\textbf{Dataset} & \shortstack{\textbf{Length}\\$\boldsymbol{(N)}$} & \shortstack{\textbf{Granules}\\$\boldsymbol{(S)}$} & \textbf{Fast-DTW} & \shortstack{\textbf{Proposed}\\\textbf{10D}} & \textbf{Speedup} \\ \hline
\textbf{GunPoint} & 150 & 6 & 167.46s & \textbf{11.78s} & \textbf{14.2$\times$} \\
\textbf{Coffee}   & 286 & 11 & 35.33s & \textbf{3.69s} & \textbf{9.5$\times$} \\
\textbf{ArrowHead}& 251 & 11 & 247.97s & \textbf{27.75s} & \textbf{8.9$\times$} \\
\textbf{ECG200}   & 96  & 5  & 42.15s  & \textbf{6.84s} & \textbf{6.2$\times$} \\ \hline
\end{tabular*}
\end{table}

\subsection{Single-Split Benchmark and Statistical Significance}
To evaluate classification performance under standard benchmark conditions, Table \ref{tab:single_split_benchmark} presents the single-split test results across all 20 UCR datasets using official train/test partitions. We compare the Proposed Adaptive LFIG (10D) model against reproduced DTW-1NN, ROCKET, and MiniROCKET baselines \cite{dempster2021}, alongside published results for HIVE-COTE 2.0 \cite{middlehurst2021}, reporting the performance difference ($\Delta = \text{Proposed} - \text{Baseline}$) for each comparison.

\begin{table*}[!t]
\caption{Single-Split Benchmark: Proposed 10D Model vs. Baselines with Performance Differences}
\label{tab:single_split_benchmark}
\centering
\ieeetableformat
\setlength{\tabcolsep}{1.5pt}
\begin{tabular}{|>{\raggedright\arraybackslash}p{3.6cm}|c|c|c|c|c|c|c|>{\raggedright\arraybackslash}p{3.1cm}|}
\hline
\multicolumn{1}{|c|}{\textbf{Dataset}} & \shortstack{\textbf{Proposed}\\\textbf{(10D)}} & \textbf{DTW-1NN} & $\Delta_{\text{DTW}}$ & \textbf{ROCKET} & $\Delta_{\text{RCK}}$ & \textbf{MiniROCKET} & $\Delta_{\text{MRCK}}$ & \multicolumn{1}{c|}{\shortstack{\textbf{Best}\\\textbf{Baseline}}} \\ \hline
GunPoint               & 0.8533 & 0.9067 & $-0.0534$ & 1.0000 & $-0.1467$ & 0.9933 & $-0.1400$ & ROCKET (1.0000) \\
Coffee                 & 0.9286 & 1.0000 & $-0.0714$ & 1.0000 & $-0.0714$ & 1.0000 & $-0.0714$ & DTW-1NN (1.0000) \\
ArrowHead              & 0.7086 & 0.7029 & \textbf{$+0.0057$} & 0.8057 & $-0.0971$ & 0.8514 & $-0.1428$ & HIVE-COTE 2.0 (0.8710) \\
ECG200                 & 0.7200 & 0.7700 & $-0.0500$ & 0.9200 & $-0.2000$ & 0.9100 & $-0.1900$ & ROCKET (0.9200) \\
Chinatown              & 0.9184 & 0.9738 & $-0.0554$ & 0.9825 & $-0.0641$ & 0.9825 & $-0.0641$ & HIVE-COTE 2.0 (0.9830) \\
ItalyPowerDemand       & 0.8863 & 0.9504 & $-0.0641$ & 0.9699 & $-0.0836$ & 0.9640 & $-0.0777$ & HIVE-COTE 2.0 (0.9700) \\
SonyAIBORobotSurface1  & 0.7837 & 0.7255 & \textbf{$+0.0582$} & 0.9168 & $-0.1331$ & 0.8935 & $-0.1098$ & ROCKET (0.9168) \\
TwoLeadECG             & 0.6778 & 0.9043 & $-0.2265$ & 0.9991 & $-0.3213$ & 0.9982 & $-0.3204$ & HIVE-COTE 2.0 (1.0000) \\
ECGFiveDays            & 0.7793 & 0.7677 & \textbf{$+0.0116$} & 1.0000 & $-0.2207$ & 1.0000 & $-0.2207$ & ROCKET (1.0000) \\
MoteStrain             & 0.8602 & 0.8347 & \textbf{$+0.0255$} & 0.9153 & $-0.0551$ & 0.9257 & $-0.0655$ & MiniROCKET (0.9257) \\
Beef                   & 0.4000 & 0.6333 & $-0.2333$ & 0.8000 & $-0.4000$ & 0.8333 & $-0.4333$ & MiniROCKET (0.8333) \\
OliveOil               & 0.8333 & 0.8333 & \textbf{$+0.0000$} & 0.9333 & $-0.1000$ & 0.9333 & $-0.1000$ & ROCKET (0.9333) \\
Meat                   & 0.6833 & 0.9333 & $-0.2500$ & 0.9500 & $-0.2667$ & 0.9667 & $-0.2834$ & MiniROCKET (0.9667) \\
BeetleFly              & 0.6000 & 0.7000 & $-0.1000$ & 0.9000 & $-0.3000$ & 0.9000 & $-0.3000$ & ROCKET (0.9000) \\
BirdChicken            & 0.7000 & 0.7500 & $-0.0500$ & 0.9000 & $-0.2000$ & 0.9000 & $-0.2000$ & ROCKET (0.9000) \\
FaceFour               & 0.7841 & 0.8295 & $-0.0454$ & 0.9773 & $-0.1932$ & 0.9886 & $-0.2045$ & MiniROCKET (0.9886) \\
SyntheticControl       & 0.9733 & 0.9933 & $-0.0200$ & 1.0000 & $-0.0267$ & 0.9833 & $-0.0100$ & ROCKET (1.0000) \\
CBF                    & 0.9000 & 0.9967 & $-0.0967$ & 1.0000 & $-0.1000$ & 0.9989 & $-0.0989$ & ROCKET (1.0000) \\
TwoPatterns            & 0.7775 & 1.0000 & $-0.2225$ & 1.0000 & $-0.2225$ & 0.9962 & $-0.2187$ & DTW-1NN (1.0000) \\
Wafer                  & 0.9455 & 0.9799 & $-0.0344$ & 0.9985 & $-0.0530$ & 0.9994 & $-0.0539$ & MiniROCKET (0.9994) \\ \hline
\end{tabular}
\end{table*}

The proposed model shows competitive performance against standard DTW-1NN, matching or exceeding its accuracy on five datasets: \texttt{SonyAIBORobotSurface1} ($+0.0582$), \texttt{MoteStrain} ($+0.0255$), \texttt{ECGFiveDays} ($+0.0116$), \texttt{ArrowHead} ($+0.0057$), and \texttt{OliveOil} ($+0.0000$), while running up to $14.2\times$ faster. While deep convolutional models (ROCKET, MiniROCKET) and ensemble architectures (HIVE-COTE 2.0) achieve higher accuracy on these static partitions due to substantially greater parameter capacity, our framework provides a competitive balance of accuracy, speed, and interpretability.

We conducted a Friedman test followed by a Nemenyi post-hoc analysis ($\alpha = 0.05$) \cite{demsar2006} across the 20 datasets. The test yielded a Friedman statistic of 45.83 ($p < 10^{-6}$), confirming statistically significant performance differences among the models. The average ranks (lower is better) are: ROCKET (1.5000), MiniROCKET (1.6750), DTW-1NN (3.0500), and Proposed Adaptive LFIG (3.7750), with a critical difference of $\text{CD} = 1.0488$. As shown in the Critical Difference diagram (Fig. \ref{fig:cd_diagram}), the proposed model and DTW-1NN fall within the same critical difference interval.

\begin{figure}[!t]
\centering
\includegraphics[width=\linewidth]{plots/cd_diagram.png}
\caption{Dem\v{s}ar Critical Difference diagram comparing average ranks of the Proposed Adaptive LFIG model, ROCKET, MiniROCKET, and DTW-1NN baselines across 20 UCR datasets.}
\label{fig:cd_diagram}
\end{figure}

\subsection{Nested Cross-Validation and Generalization Performance}
A significant methodological issue in time series benchmarks is selection leakage, where hyperparameter choices are inadvertently influenced by test set feedback. To quantify this effect, we compared test-guided parameter selection with our leakage-free inner cross-validation protocol. When hyperparameters (spread factor $z$, neighbor count $k$, and fusion weights $\mathbf{w}$) were tuned using test set feedback, accuracy was inflated by 16.00\% on \texttt{ECG200} (0.8800 vs. 0.7200) and by 5.34\% on \texttt{GunPoint} (0.9067 vs. 0.8533), whereas stationary datasets such as \texttt{Coffee} remained stable (0.9286). This confirms that signals with notable phase variability are particularly vulnerable to over-optimistic test evaluation.

\begin{table*}[!t]
\caption{5-Fold Nested Cross-Validation: Generalization Results, Macro F1, and Routing Strategy}
\label{tab:nested_cv_generalization}
\centering
\ieeetableformat
\setlength{\tabcolsep}{2pt}
\begin{tabular}{|>{\raggedright\arraybackslash}p{3.6cm}|>{\raggedright\arraybackslash}p{1.6cm}|c|c|c|c|}
\hline
\multicolumn{1}{|c|}{\textbf{Dataset}} & \multicolumn{1}{c|}{\textbf{Domain}} & \shortstack{\textbf{Nested CV Acc}\\\textbf{(Mean $\pm$ SD)}} & \shortstack{\textbf{Nested CV F1}\\\textbf{(Mean $\pm$ SD)}} & \shortstack{\textbf{Optimal}\\\textbf{Strategy}} & \shortstack{\textbf{Tuned Weights}\\$\boldsymbol{(w_H, w_{\text{DTW}}, w_{\text{Cos}})}$} \\ \hline
GunPoint               & Motion    & 0.9300 $\pm$ 0.0100   & 0.9300 $\pm$ 0.0100   & CPD(1.5)    & (0.2, 0.4, 0.4) \\
Coffee                 & Spectro   & 1.0000 $\pm$ 0.0000   & 1.0000 $\pm$ 0.0000   & Fixed(28)   & (0.1, 0.8, 0.1) \\
ArrowHead              & Image     & 0.8910 $\pm$ 0.0241   & 0.8917 $\pm$ 0.0226   & Fixed(25)   & (0.8, 0.1, 0.1) \\
ECG200                 & ECG       & 0.8000 $\pm$ 0.0671   & 0.7788 $\pm$ 0.0673   & CPD(1.5)    & (0.4, 0.2, 0.4) \\
Chinatown              & Traffic   & 0.9807 $\pm$ 0.0140   & 0.9770 $\pm$ 0.0165   & Fixed(10)   & (0.1, 0.8, 0.1) \\
ItalyPowerDemand       & Sensor    & 0.9644 $\pm$ 0.0159   & 0.9644 $\pm$ 0.0159   & CPD(1.5)    & (0.1, 0.1, 0.8) \\
SonyAIBORobotSurface1  & Sensor    & 0.9742 $\pm$ 0.0156   & 0.9739 $\pm$ 0.0158   & Fixed(10)   & (0.2, 0.6, 0.2) \\
TwoLeadECG             & ECG       & 0.9845 $\pm$ 0.0034   & 0.9845 $\pm$ 0.0034   & Fixed(10)   & (0.1, 0.1, 0.8) \\
ECGFiveDays            & ECG       & 0.9819 $\pm$ 0.0109   & 0.9819 $\pm$ 0.0110   & Fixed(13)   & (0.1, 0.8, 0.1) \\
MoteStrain             & Sensor    & 0.9465 $\pm$ 0.0138   & 0.9462 $\pm$ 0.0139   & Fixed(10)   & (0.1, 0.8, 0.1) \\
Beef                   & Spectro   & 0.3667 $\pm$ 0.0850   & 0.3611 $\pm$ 0.0493   & CPD(1.5)    & (0.1, 0.8, 0.1) \\
OliveOil               & Spectro   & 0.7500 $\pm$ 0.0745   & 0.6250 $\pm$ 0.0445   & CPD(1.5)    & (0.2, 0.4, 0.4) \\
Meat                   & Spectro   & 0.8917 $\pm$ 0.0333   & 0.8909 $\pm$ 0.0348   & CPD(1.5)    & (0.3, 0.4, 0.3) \\
BeetleFly              & Image     & 0.7750 $\pm$ 0.0935   & 0.7697 $\pm$ 0.0957   & Fixed(51)   & (0.4, 0.2, 0.4) \\
BirdChicken            & Image     & 0.8000 $\pm$ 0.0612   & 0.7959 $\pm$ 0.0633   & Fixed(51)   & (0.8, 0.1, 0.1) \\
FaceFour               & Image     & 0.9285 $\pm$ 0.0363   & 0.9253 $\pm$ 0.0428   & Fixed(35)   & (0.3, 0.4, 0.3) \\
SyntheticControl       & Simulated & 0.9933 $\pm$ 0.0062   & 0.9933 $\pm$ 0.0062   & Fixed(10)   & (0.4, 0.2, 0.4) \\
CBF                    & Simulated & 0.9817 $\pm$ 0.0111   & 0.9817 $\pm$ 0.0111   & CPD(1.5)    & (0.1, 0.1, 0.8) \\
TwoPatterns            & Simulated & 0.7110 $\pm$ 0.0112   & 0.7107 $\pm$ 0.0116   & CPD(1.5)    & (0.1, 0.1, 0.8) \\
Wafer                  & Sensor    & 0.9782 $\pm$ 0.0038   & 0.9393 $\pm$ 0.0113   & CPD(1.5)    & (0.1, 0.1, 0.8) \\ \hline
\end{tabular}
\end{table*}

Table \ref{tab:nested_cv_generalization} presents the multi-fold generalization performance of the proposed architecture under 5-fold outer / 3-fold inner nested cross-validation. All hyperparameters, segmentation strategies, and metric weights are chosen strictly within the inner folds. Evaluating across outer folds provides an unbiased assessment of generalization: the model achieves mean accuracies of $1.0000 \pm 0.0000$ on \texttt{Coffee}, $0.9807 \pm 0.0140$ on \texttt{Chinatown}, $0.9742 \pm 0.0156$ on \texttt{SonyAIBORobotSurface1}, and $0.9300 \pm 0.0100$ on \texttt{GunPoint}.

\section{Discussion}

\subsection{Dynamic CPD vs. Fixed-Window Partitioning}
One of the key observations from our benchmark is the structural dichotomy between phase-shifted and phase-aligned series. For phase-shifted signals (such as \texttt{GunPoint} and \texttt{ECG200}), volatile patterns exhibit regime transitions at variable time offsets. In these cases, adaptive CPD dynamically aligns window boundaries with structural transitions, reducing regression residuals. Conversely, for phase-aligned signals (such as \texttt{ArrowHead}, \texttt{Chinatown}, and \texttt{Coffee}), signals are stationary or length-normalized with synchronized events across samples. Fixed-window partitioning preserves the temporal alignment of granules across series, avoiding spurious shifts introduced by CPD.

\subsection{Feature Importance and Ablation}
Our 10-dimensional granule representation provides noticeable classification improvements under both adaptive CPD ($+2.67\%$ on GunPoint) and fixed windowing ($+2.29\%$ on ArrowHead) relative to standard 3D configurations. However, on short phase-shifted signals like ECG200, extracting 10 features over very few granules (e.g., 5 granules) can introduce noise, leading to an accuracy drop of $-8.00\%$. On other phase-aligned datasets under fixed windowing (e.g., Coffee), the additional features preserve high accuracy without degradation.

The Leave-One-Feature-Out (LOFO) ablation results across the UCR cohort are displayed in Fig. \ref{fig:lofo_importance}.

\begin{figure}[!t]
\centering
\includegraphics[width=\linewidth]{plots/fig3_lofo_importance.png}
\caption{Leave-One-Feature-Out (LOFO) ablated accuracy impact delta ($\Delta = \text{Full} - \text{Ablated}$) with 95\% confidence intervals across UCR datasets, illustrating the direct utility of primary features (positive deltas) and the stabilizing effect of global descriptors (negative deltas).}
\label{fig:lofo_importance}
\end{figure}

\subsection{Speedup Analysis}
Fuzzy granulation compresses raw time series of length $N$ into $S \ll N$ granules. Since DTW complexity scales quadratically with sequence length, computing DTW over $S$ granules rather than $N$ raw points significantly reduces alignment operations. This explains why the proposed framework achieves up to a $14.2\times$ speedup over raw Fast-DTW while maintaining competitive classification accuracy.

\subsection{Statistical Power and Wilcoxon Limitations}
With only $n=5$ datasets, the Wilcoxon signed-rank test has limited statistical power, as the $p$-value cannot fall below 0.0625 regardless of effect size. Expanding the evaluation cohort to 20 datasets allows using the Friedman test with Nemenyi post-hoc analysis, providing a more reliable assessment of performance differences across models.

\subsection{Interpretability and Computational Trade-Offs}
Selecting a time series classification model often involves trade-offs among classification accuracy, computational efficiency, memory consumption, and explainability. Recent state-of-the-art methods achieve high accuracy using deep ensemble architectures (e.g., HIVE-COTE 2.0) or extensive random convolutional transforms (e.g., ROCKET and MiniROCKET). However, these approaches carry practical drawbacks: HIVE-COTE requires significant training time and memory, while ROCKET projects signals into tens of thousands of random features whose individual contributions cannot be readily inspected.

In contrast, the proposed Adaptive Multi-Feature LFIG framework provides a practical alternative suited for resource-constrained or interpretability-sensitive applications:
\begin{enumerate}
    \item \textit{Computational Efficiency:} By compressing a raw series of length $N$ into $S \ll N$ fuzzy granules, dynamic programming operates in $O(S^2)$ rather than $O(N^2)$, achieving up to a $14.2\times$ speedup over DTW and reducing peak heap memory by up to $5.5\times$. This makes the model well-suited for edge devices and microcontroller platforms.
    \item \textit{Model Interpretability:} Unlike high-dimensional black-box embeddings, each descriptor in our 10D representation corresponds to a concrete signal property (e.g., trend slope, Shannon entropy, volatility, curvature, and energy). These features allow domain experts to examine decision rationales directly alongside the fuzzy envelope bounds $[L, U]$. For example, in ECG monitoring, clinicians can inspect whether an arrhythmia detection was driven by local slope changes, energy spikes, or entropy shifts.
    \item \textit{Analytical Smoothing:} Fitting least-squares trends within segments serves as an intrinsic filter against high-frequency sensor noise while preserving the envelope bounds to reflect measurement uncertainty.
\end{enumerate}

\section{Conclusion}
In this paper, we proposed an Adaptive Multi-Feature Linear Fuzzy Information Granulation framework for time series classification. By dynamically windowing signals, extracting 10-dimensional statistical-structural feature spaces, and fusing Hausdorff, slope-DTW, and cosine-DTW metrics, our framework achieves classification accuracy competitive with standard DTW baselines while running up to $14.2\times$ faster. Although large ensemble models (such as HIVE-COTE 2.0) and convolutional feature extractors (such as ROCKET) achieve higher accuracy on official single splits, our approach delivers substantial savings in computation and memory along with transparent, human-interpretable granule representations.

\subsection{Future Work}
Future work will explore three directions. First, we plan to extend the multi-feature granular representation and hybrid metrics to multivariate time series, exploring inter-channel co-granulation mechanisms. Second, to accelerate pairwise distance computations on larger datasets, we intend to implement GPU-accelerated dynamic programming routines. Third, we plan to incorporate prototype selection techniques, such as $k$-medoids or condensed nearest neighbors, to reduce the number of reference templates and further lower inference latency.

\begin{thebibliography}{10}

\bibitem{gao2019}
F. Gao and F. Yu, ``Linear fuzzy information granulation based classification method for unequal length time series,'' \textit{IEEE Access}, vol. 7, pp. 91118--91128, 2019.

\bibitem{he2023}
Q. He and F. Yu, ``Trend recurrence analysis and time series classification via trend fuzzy granular recurrence plot method (TFGRP-SVM),'' \textit{Chaos, Solitons \& Fractals}, vol. 169, p. 113309, 2023.

\bibitem{duan2018}
L. Duan, F. Yu, W. Pedrycz, and Y. Wang, ``Time-series clustering based on linear fuzzy information granules,'' \textit{Applied Soft Computing}, vol. 73, pp. 1053--1067, 2018.

\bibitem{yang2017}
X. Yang, F. Yu, and W. Pedrycz, ``Long-term forecasting of time series based on linear fuzzy information granules and fuzzy inference system,'' \textit{International Journal of Approximate Reasoning}, vol. 81, pp. 1--27, 2017.

\bibitem{demsar2006}
J. Dem\v{s}ar, ``Statistical comparisons of classifiers over multiple data sets,'' \textit{Journal of Machine Learning Research}, vol. 7, pp. 1--30, 2006.

\bibitem{dempster2021}
A. Dempster, F. Schmidt, and G. I. Webb, ``MINIROCKET: A Very Fast and Accurate Language for Time Series Classification,'' in \textit{Proceedings of the 27th ACM SIGKDD Conference on Knowledge Discovery \& Data Mining}, 2021, pp. 248--257.

\bibitem{middlehurst2021}
M. Middlehurst, J. Large, M. Flynn, J. Featherstone, and J. Bagnall, ``HIVE-COTE 2.0: a new meta-ensemble for time series classification,'' \textit{Machine Learning}, vol. 110, no. 11, pp. 3211--3243, 2021.

\bibitem{truong2020}
C. Truong, L. Oudre, and N. Vayatis, ``Selective review of offline change point detection methods,'' \textit{Signal Processing}, vol. 167, p. 107299, 2020.

\bibitem{sakoe1978}
H. Sakoe and S. Chiba, ``Dynamic programming algorithm optimization for spoken word recognition,'' \textit{IEEE Transactions on Acoustics, Speech, and Signal Processing}, vol. 26, no. 1, pp. 43--49, 1978.

\bibitem{lubba2019}
C. H. Lubba, S. S. Sethi, P. Knaute, S. R. Schultz, B. D. Fulcher, and N. S. Jones, ``catch22: CAnonic Time-series CHaracteristics on 22 non-redundant features,'' \textit{Data Mining and Knowledge Discovery}, vol. 33, no. 6, pp. 1823--1846, 2019.

\end{thebibliography}

\end{document}
